%%%%%%%%%%%%%%%%%%%%%%%%
%
% $Autor: Wings $
% $Datum: 2020-07-24 09:05:07Z $
% $Pfad: GDV/Vortraege/latex - Ausarbeitung/Kapitel/Bezier.tex $
% $Version: 4732 $
%
%%%%%%%%%%%%%%%%%%%%%%%%

\chapter{Power source}
%Elektor Special: Stromversorgung & Batterien  

All the Arduino boards need power to operate, either it comes from the USB connection with Laptop, Ac power adapter, Battery or a regulated power supply. \\
the USB connection is the easiest way for simple testing The Arduino Nano 33 BLE Sense can be powered through a USB cable connected to a computer or USB power supply. This is the most common way to power the device during development and testing.
\section{Power sources Available}

(introduction, compatibility with our board, how long does it last, advantages disadvantages)
External power supply: The Arduino Nano 33 BLE Sense can also be powered by an external power supply, such as a battery or a power adapter. The device can be powered by a voltage between 7 and 12 V, and it has a built-in voltage regulator that converts the input voltage to the 3.3V required by the device.

VIN pin: The Arduino Nano 33 BLE Sense also has a VIN pin, which can be used to supply power to the device. This pin is connected to the voltage regulator, and it can be used to power the device with an external voltage source.

Li-Po battery: The Arduino Nano 33 BLE Sense also has a Li-Po battery connector, which can be used to power the device with a Li-Po battery.
\section{Choice+ how to}

